\documentclass[graybox]{svmult}

% ---------------------------------------------------------
% Packages
% ---------------------------------------------------------
\usepackage[T1]{fontenc}
\usepackage{type1cm}
\usepackage{graphicx}
\usepackage{multicol}
\usepackage[bottom]{footmisc}
\usepackage{amsmath,amssymb}
\usepackage{url}
\usepackage{booktabs}
\usepackage{multirow}
\usepackage{array}
\usepackage[type1]{roboto}
\usepackage{algorithm}
\usepackage{algpseudocode}
\usepackage{float}

% Times-like text and mathematics, matching the supplied Springer template.
\IfFileExists{newtxtext.sty}{%
  \usepackage{newtxtext}
  \usepackage[varvw]{newtxmath}
}{%
  \usepackage{mathptmx}
}

\definecolor{navyblue}{RGB}{0,51,102}

% A non-floating algorithm environment that may continue across pages.
\newenvironment{breakablealgorithm}
{%
    \par\medskip
    \refstepcounter{algorithm}%
    \hrule\vspace{2pt}%
    \renewcommand{\caption}[1]{%
        \noindent\textbf{Algorithm~\thealgorithm}\enspace ##1\par
        \vspace{2pt}\hrule\vspace{2pt}%
    }%
}
{%
    \vspace{2pt}\hrule
    \par\medskip
}


% ---------------------------------------------------------
% Document
% ---------------------------------------------------------
\begin{document}

% ---------------------------------------------------------
% Title
% ---------------------------------------------------------
\title*{An Emotion-Aware Conversational AI System with Tiered and
Multi-Label Affective Memory}
\titlerunning{Emotion-Aware AI with Tiered and Multi-Label Affective Memory}

\author{Mayank Kumar \and Sujeet Kumar Sah \and Sagnik Sen}
\authorrunning{Mayank Kumar et al.}

\institute{
Mayank Kumar \at Department of Computer Science \& Engineering,
Institute of Engineering \& Management, Salt Lake, Kolkata, India,
\email{mayank.kumar2023@iem.edu.in}
\and
Sujeet Kumar Sah \at Department of CSE (AI),
Institute of Engineering \& Management, Salt Lake, Kolkata, India,
\email{sujeet.kumarsah2023@iem.edu.in}
\and
Sagnik Sen \at Department of Computer Science \& Engineering,
Institute of Engineering \& Management, Salt Lake, Kolkata, India,
\email{Sagnik.Sen2023@iem.edu.in}
}

\maketitle

% ---------------------------------------------------------
% Abstract
% ---------------------------------------------------------
\abstract{
LLM conversational AI systems  fails to analyze convoluted emotions at different sessions  and often treat Emotion and Memory seperately.  We introduce a unified emotion-aware architecture integrating parameter-efficient multi-label emotion recognition, tiered affective memory, and hybrid graph-vector retrieval. A
parameter-efficient model fine-tuned via Low-Rank Adaptation on the 28-class GoEmotions dataset employs sigmoid-based multi-label classification to detect co-occurring emotions. A three-tier affective memory—situational, short-term, and long-term—with separate half-life decays and
a Valence-Arousal-Dominance reasoning core making it possible to capture emotional patterns and achieve 90\% of relevant signals in a
relatively short time (2, 8, and 21 turns accordingly). ContextManager orchestrates Redis caching, Neo4j graph memory, and vector-based Retrieval-Augmented
Generation to apply smaller-scale
affective summaries while keeping in mind all token limitations. An empathetic response-composition layer recognizes concern without aggravating it, while a deterministic safety override handles crisis inputs.Controlled validations confirm decay correctness (semigroup error $\le 2.5\times10^{-16}$), weak-signal
label recovery, ($0\% \rightarrow 66.7\%$), 100\% compliance with
the token budget for 100-20,000 memories, and full episodic recall during history eviction. Code is available at {\roboto\mdseries\urlstyle{same}\textcolor{navyblue}{\url{https://github.com/sujeet-sah-748/Emotion-Aware-Conversation-AI-System/}}}}

% ---------------------------------------------------------
% Keywords
% ---------------------------------------------------------
\keywords{Tiered Memory Architecture, VAD emotional Model, Temporal affect decay, Affective Memory, Crisis Aware Responses, Emotional Trajectory Modeling, Parameter Efficient Fine-tuning, Hybrid Memory Systems}


% =========================================================
% I. INTRODUCTION
% =========================================================

\section{Introduction}
\label{sec:introduction}

Conversational AI research has progressed from stateless language generation to systems in which LLMs are augmented with retrieved information, hierarchical memory, and experience~\cite{lewis2020rag,packer2023memgpt,park2023generative}. Generating emotion-aware responses involves more than semantic comprehension. Traditional sentiment analysis collapses the entire utterance into either positive, negative, or neutral sentiments, discarding the nuance of natural dialogue, where a single utterance often conveys multiple, sometimes conflicting, emotions (e.g., “I am terrified but excited about tomorrow” expresses both fear and anticipation).

Another problem that affects long-term conversations is that of the non-persistence of emotional context in a chat beyond sessions. While contextual emotion detection systems such as DialogueRNN and COSMIC have been proposed for the detection of emotion evolution within a fixed context window, they do not address the issue of how this state should persist, decay, or be recalled when the dialogue continues beyond sessions~\cite{majumder2019dialoguernn,ghosal2020cosmic}. In contrast, architectures for long-horizon memory retention, such as MemGPT, Generative Agents, and LongMemEval~\cite{packer2023memgpt,park2023generative,wu2025longmemeval}, operate with the retention of facts and events but not with the development or degradation of emotional states of the user. Hence, after the emotions-relevant turn has been summarized or truncated to the extent of fitting the bounded context window, the affective information is lost either completely or gets mixed up with other pieces of information, and there is no way of marking it as relevant anymore.

In this paper, we propose a unifying emotional architecture which takes into consideration all factors of using emotional knowledge efficiently in a dialogue: multi-label emotion recognition, three-level emotional memory, hybrid retrieval, and generation of responses taking crisis issues into account while using the context information and the token budget for building the context of the response. Regarding the perception component of the proposed system, the DistilRoBERTa which is trained using the LoRA approach~\cite{hu2021lora} will be used for classifying the data into 28 categories given by the GoEmotions dataset~\cite{demszky2020goemotions} with the help of sigmoid multi-label classification. As for the memory component of the system, we present a three-level emotional memory structure, which is capable of storing information in situational (short-term) and long-term fashion and which proposes the concept of half-life decay independently for the emotion storage process and, despite that, allows for the use of VAD reasoning core. These states are stored and accessed by means of a hybrid stack of Redis caching, vector-based RAG, and Neo4j graph memory. At the generation component, an empathic-aware response composition module operates according to the concepts of empathic and emotionally supportive conversation~\cite{rashkin2019empathetic,liu2021esconv}, while a deterministic safety override takes care of crisis-related inputs. \linebreak

\textbf{Our Contribution:} 

\begin{enumerate}


\item \textbf{Tiered affective memory with decay:} A three-tier (situational, short-term, long-term) architecture with independent half-life decays, a VAD reasoning core, and trajectory modeling of multi-label affect over time.

\item \textbf{Dual-representation design (VAD reasoning core + multi-label emotion bars):}
A decoupled architecture maintaining (i) a continuous VAD vector for decay computation, threshold logic, and event detection, and (ii) parallel 28-label probability bars for interpretable display and retrieval, preventing information loss while enabling both numerical reasoning and human-readable emotion tracking.

\item \textbf{Hybrid graph–vector memory integration:} Unification of Redis, vector-based RAG, and Neo4j graph memory within a token-budgeted pipeline that injects a size-bounded live affect summary into every prompt.

\item \textbf{Safety-aware generation:} A response-composition strategy that validates negative emotions before broadening perspective, coupled with a hard override for crisis-related inputs.

\item \textbf{Concrete defect analysis:} Identification and correction of three implementation-level issues—a magnitude-biased nearest-neighbor label projector, a configuration mismatch in session-reset paths, and a multi-label-to-single-vector collapse—with controlled quantification of their impact.

\end{enumerate}




% =========================================================
% II. RELATED WORK
% =========================================================
\section{Literature Review}

Recent developments in conversational artificial intelligence have shifted the focus from task-oriented and stateless chatbots toward systems capable of maintaining conversational context, understanding user emotions, and providing personalized responses~\cite{zhong2024memorybank,packer2023memgpt,rashkin2019empathetic,liu2021esconv}. Existing research has explored these capabilities through memory-enhanced architectures, emotion recognition models, context-aware conversational systems, and human-centered interaction studies. However, these capabilities are often investigated independently, creating an opportunity for an integrated personalized conversational framework.

\subsection{Memory-Enhanced and Personalized Conversational Systems}

Memory is an important component of personalized conversational AI because conventional chatbots generally lose information when a conversation ends. Fathi et al.~\cite{Fathi2026} proposed an adaptive memory system for personalized chatbots that combines short-term and long-term memory. The short-term memory maintains the immediate conversational context, while long-term memory preserves user-specific information across sessions. Their framework uses machine learning and natural language processing techniques to dynamically store and retrieve relevant information. The study evaluates response accuracy, user satisfaction, interaction fluidity, memory retrieval performance, and computational efficiency. The results indicate that adaptive memory can improve personalization, continuity, and user engagement. However, the increasing size of long-term memory introduces challenges related to computational cost, memory relevance, privacy, and the selection of information that should be retained.

Oni~\cite{Oni2025} developed a memory-enhanced generative conversational AI system designed to improve context awareness and personalization. The system employs a generative conversational architecture with memory storage and retrieval to recall previous interactions and generate contextually relevant responses. The study evaluates fluency, user satisfaction, memory recall, perplexity, diversity, and consistency. The findings demonstrate that memory retrieval can reduce repetitive interactions and improve the continuity of conversations. Nevertheless, the study primarily focuses on conversational memory and does not explicitly model the emotional state of the user when selecting or updating memories.

Jover~\cite{Jover2026} presented AILA, an empathic voice assistant incorporating persistent memory and emotion awareness for elderly users with cognitive decline. The system uses Python and FastAPI for the backend, Mem0 and ChromaDB for memory management, and Hume AI for voice interaction and emotion detection. A two-layer memory architecture is used to distinguish structured user-profile information from conversational memories. Iterative functional testing showed persistent memory across sessions, hands-free interaction, and successful integration of emotion-aware functionality. Although the system demonstrates the practical feasibility of combining memory and emotion awareness, its evaluation is primarily prototype-oriented and does not provide extensive independent validation of emotion recognition or long-term user outcomes.

Jiang et al.~\cite{Jiang2026} introduced RECALLbot, an LLM-based social chatbot that extends conventional user memory with agentic and reciprocal memory. The system introduces ``Me Memory,'' representing the chatbot's own experiences, and ``We Memory,'' representing shared experiences between the user and chatbot. A two-week study involving 40 participants demonstrated improvements in perceived social identity, self-disclosure, and trust compared with a baseline chatbot. This work demonstrates that memory can influence not only response personalization but also the quality of human-chatbot relationships. However, the relatively small sample size and short evaluation period limit the generalizability of the findings.

Cox et al.~\cite{Cox2026} further investigated the relationship between chatbot memory and user self-disclosure. Their study examined how persistent memory and conversational framing influence users' decisions to disclose personal information. The study highlights an important dual effect of memory: persistent memory can increase continuity and familiarity, but it can also raise privacy concerns and influence users' perceptions of how their information is being retained and used. This finding emphasizes the need for memory transparency, user control, and privacy-aware memory mechanisms in personalized conversational systems.

Related LLM-memory architectures provide complementary mechanisms. MemoryBank stores, recalls, and updates long-term user information to support continued interaction~\cite{zhong2024memorybank}; MemGPT manages limited context through a hierarchy inspired by operating-system memory~\cite{packer2023memgpt}; and Generative Agents retrieve memories using relevance, recency, and importance~\cite{park2023generative}. More recent approaches organize memory through graph structures or agent-driven updates, including GraphRAG, HippoRAG, A-MEM, Zep, and Mem0~\cite{edge2024graphrag,gutierrez2024hipporag,xu2025amem,rasmussen2025zep,chhikara2025mem0}. LongMemEval and long-horizon conversational-memory studies further show that retaining and using information across extended interactions remains a distinct evaluation challenge~\cite{wu2025longmemeval,maharana2024longterm}.

\subsection{Emotion-Aware Conversational Systems}

In addition to memory, emotion recognition has become an important research direction for improving the quality of conversational interaction. Akram et al.~\cite{Akram2026} conducted a systematic review of 50 studies related to emotion detection in chatbot systems. The review examined emotion recognition techniques, data modalities, architectures, and evaluation methodologies. The authors observed that text-based emotion detection remains highly scalable, whereas speech, visual, and multimodal approaches provide richer emotional information. The review also identifies a significant research gap: many existing systems emphasize technical performance metrics such as classification accuracy while giving comparatively less attention to long-term human-centered outcomes such as trust, engagement, continuity, and perceived intelligence. The authors therefore emphasize the need for ethically responsible, culturally sensitive, and systematically evaluated emotion-aware conversational systems.

Abinaya et al.~\cite{Abinaya2025} proposed TEBC-Net, a multimodal emotion-aware conversational agent that combines BERT-based text emotion recognition with CNN-based facial emotion recognition. The text component achieved an accuracy of 87.21\%, while the visual component achieved 74.14\%. Facial preprocessing techniques, including Haar Cascade detection, CLAHE, and dual gamma correction, were employed to improve visual emotion recognition. The outputs from text and facial analysis were subsequently used to generate more empathetic chatbot responses. Although the multimodal design provides richer emotional information, the approach introduces additional computational complexity and privacy considerations. Moreover, the system depends on visual input, which may not be appropriate for a text-centered personalized conversational environment.

Devi et al.~\cite{Devi2026} proposed a multimodal context-aware conversational agent for emotion recognition and adaptive stress management. The system integrates textual, acoustic, and visual information using deep learning models, including CNN and LSTM architectures. The reported multimodal emotion recognition accuracy was 93.6\%, outperforming CNN-only, LSTM-based, and transformer-based approaches. Contextual intelligence also improved response relevance, while personalized stress-management strategies achieved high user satisfaction. The study demonstrates the value of combining emotional information with contextual information. However, its multimodal design is primarily oriented toward stress-management interventions and requires additional consideration of long-term personalized memory and privacy-aware conversational adaptation.

Text-based dialogue research has also established benchmarks for empathetic response generation and emotional support. EmpatheticDialogues grounds conversations in emotional situations and evaluates perceived empathy~\cite{rashkin2019empathetic}, whereas ESConv explicitly annotates emotional-support strategies~\cite{liu2021esconv}. For contextual emotion recognition, DialogueRNN models speaker-aware conversational dynamics~\cite{majumder2019dialoguernn}; EmotionLines and MELD provide multi-party conversational emotion corpora~\cite{chen2018emotionlines,poria2019meld}; and COSMIC incorporates commonsense knowledge into emotion identification in conversation~\cite{ghosal2020cosmic}. These studies demonstrate the importance of conversational context, but they do not by themselves provide a persistent, user-controlled affective memory across sessions.

\subsection{Integrated Emotion, Context, Memory, and Safety}

Raghav et al.~\cite{Raghav2026} proposed an emotionally intelligent conversational agent called Sentient that integrates emotion recognition, long-term memory, retrieval-enhanced reasoning, and an ethical interaction layer. The architecture attempts to address the limitations of conventional task-oriented conversational systems by combining knowledge retrieval with emotional and historical context. The proposed framework demonstrates an important direction toward integrating cognition, emotion, personalization, and ethical response generation within a single architecture. However, the use of retrieval-enhanced knowledge makes the system broader than a purely personalized conversational agent whose primary source of information is the user's conversational history and preferences. Furthermore, the integration of several components increases system complexity.

Chang and Herath~\cite{ChangHerath2025} proposed an architecture for embodied AI that combines persistent memory, emotional intelligence, attention, learning, and proactive engagement. Their work emphasizes the importance of long-term interaction in developing meaningful human-AI relationships. Persistent memory allows the system to recall previous interactions, while emotional intelligence enables the system to adapt its responses to the user's affective state. The architecture provides a broader perspective on relationship-oriented AI, although the work is primarily architectural and conceptual rather than a large-scale quantitative evaluation of personalized conversational performance.

\subsection{Comparison of Existing Approaches}

The reviewed studies demonstrate that memory and emotion are two major components for developing personalized conversational AI. Memory-based approaches primarily focus on continuity, personalization, and retrieval of previous user information~\cite{zhong2024memorybank,packer2023memgpt,park2023generative}, whereas emotion-aware approaches focus on identifying the user's affective state and adapting response generation~\cite{rashkin2019empathetic,liu2021esconv,majumder2019dialoguernn}. Some recent systems attempt to combine these capabilities with knowledge retrieval and safety mechanisms~\cite{Raghav2026,Jover2026}.

However, several limitations remain. First, representative memory systems retrieve information according to semantic relevance, recency, importance, or learned graph structure without explicitly modeling the emotional significance of each memory~\cite{park2023generative,zhong2024memorybank,edge2024graphrag,gutierrez2024hipporag}. Second, emotion recognition is frequently evaluated as classification over individual utterances or bounded conversations rather than as a persistent user state~\cite{demszky2020goemotions,majumder2019dialoguernn,ghosal2020cosmic}. Third, many systems use multimodal inputs such as facial expressions and speech, which increase sensing requirements and raise additional privacy considerations~\cite{Abinaya2025,Devi2026,poria2019meld}. Fourth, long-horizon studies show that retaining and correctly using information over extended interactions remains difficult~\cite{wu2025longmemeval,maharana2024longterm}. Finally, persistent memory can alter users' disclosure decisions, reinforcing the need for transparent retention and user control~\cite{Cox2026,chhikara2025mem0}.

\subsection{Research Gap}

The literature review indicates a clear opportunity for a personalized conversational architecture that integrates \textbf{emotion awareness, conversational context, long-term memory, and safety-aware response generation} without relying on multimodal inputs or external knowledge retrieval as its primary mechanism.

Existing memory-based systems demonstrate the benefits of persistent user information but generally do not model emotional states as part of memory formation and retrieval~\cite{zhong2024memorybank,packer2023memgpt,park2023generative,xu2025amem}. In contrast, emotion-aware systems identify user affect and support empathetic responses but often lack mechanisms for retaining and reusing emotionally relevant information across sessions~\cite{rashkin2019empathetic,liu2021esconv,majumder2019dialoguernn}. Furthermore, existing integrated architectures may depend on multimodal signals or retrieval-augmented knowledge sources~\cite{Abinaya2025,Devi2026,Raghav2026}, which are not necessary for a system whose primary objective is personalized conversation based on the user's interaction history.

Therefore, the proposed research focuses on an integrated text-based architecture in which the user's current conversational context, emotional state, preferences, and relevant long-term memories are jointly used during response generation. The proposed system further introduces an emotion-aware memory mechanism in which emotional information can be represented continuously, updated over time, and selectively persisted according to confidence, relevance, recency, and importance. A safety and privacy layer is incorporated before the final response is delivered to ensure that personalization does not result in inappropriate disclosure or unsafe responses.

This approach aims to bridge the gap between conventional memory-enhanced chatbots and emotion-aware conversational agents by treating \textbf{emotion, context, and long-term personalized memory as interconnected components of conversational intelligence}.


% =========================================================
% III. METHODOLOGY
% =========================================================

\section{Proposed Method}
\label{sec:methodology}

We propose an affect-aware context-management architecture for conversational language models. The architecture combines two complementary mechanisms: (i) a tiered affect tracker that converts per-turn emotion predictions into temporally persistent emotional states and (ii) a long-term memory and context engine that retrieves relevant memories and combines them with recent dialogue and affective information under an explicit token budget. A context manager orchestrates these mechanisms. For each user turn, the processing sequence comprises message insertion, affect inference, affect-state updating, optional emotional-memory writing, semantic memory retrieval, affect-state rendering, and token-budgeted context assembly.

The implementation deliberately separates affect \emph{reasoning} from affect \emph{presentation}. The VAD representation (valence, arousal, and dominance), grounded in established dimensional models of affect~\cite{russell1980circumplex,mehrabian1996pad}, provides the numerical state used for thresholding, temporal decay, event detection, and baseline adaptation. In parallel, a multi-label representation preserves the raw classifier probabilities as interpretable emotion bars. This separation prevents the loss of co-occurring emotions when a multidimensional affective state is projected onto a single categorical label.


Figure~\ref{fig:tiered-memory-architecture} summarizes the complete
processing and memory flow.

\begin{figure}[!t]
    \centering
    \includegraphics[width=0.95\textwidth]{architecture.jpg}
    \caption{Tiered memory architecture in which three temporal layers
    (situational, short-term, and long-term) process user interactions after
    DistilRoBERTa classification. The system maintains VAD affective states
    and behavioral metrics across a 12-pair rolling window and uses vector
    databases for context retrieval and safety-filtered response generation.}
    \label{fig:tiered-memory-architecture}
\end{figure}



\begin{breakablealgorithm}
\small
\caption{End-to-End Affective Context Processing and Memory Retrieval}
\label{alg:pipeline}
\begin{algorithmic}[1]

    \Require User utterance $u_t$, history $H$, emotion prototypes $\{q_l\}$,
    and tiered affect states $\{s^{\text{tier}}, b^{\text{tier}},
    h^{\text{tier}}, \eta^{\text{tier}}\}$
    \Statex \qquad for $\text{tier} \in \{\text{sit}, \text{short},
    \text{long}\}$; rolling window $W$; memory store $\mathcal{M}$
    \Statex \qquad token budget $B$; top-$n$ confidence width $n$; and
    thresholds $\{\gamma_{\text{sit}}, \gamma_{\text{short}},
    \gamma_{\text{long}}, \tau, \theta_{\text{spike}}, \alpha\}$

    \State Append utterance: $H \gets H \cup \{u_t\}$
    \State Extract recent context $x$ from $H$; compute multi-label probabilities $p_{t,l} = \sigma(f_\theta(x)_l)$
    \State Aggregate VAD vector $\tilde{e}_t \gets \frac{\sum_l p_{t,l}\, q_l}{\sqrt{\max\!\left(1,\, \sum_l p_{t,l}\right)}}$
    \State Compute $\bar{p}_{t,n} \gets \frac{1}{n}\sum_{k=1}^{n}p_{t,(k)}$ and $c_t \gets \text{clip}_{[0,1]}\!\left(\frac{\bar{p}_{t,n}+(\bar{p}_{t,n}-p_{t,(n+1)})}{2}\right)$

    \If{$c_t < \gamma_{\text{sit}}$}
        \State Skip affect updates and proceed directly to retrieval (line~\ref{line:retrieval})
    \EndIf

    \State Compute decayed situational state $s^{\text{eff}}_{\text{sit}} \gets b_{\text{sit}} + (s_{\text{sit}} - b_{\text{sit}})\, 2^{-\Delta t / h_{\text{sit}}}$
    \State Update situational tier: $s_{\text{sit}} \gets s^{\text{eff}}_{\text{sit}} + \eta_{\text{sit}}\, c_t\, (\tilde{e}_t - s^{\text{eff}}_{\text{sit}})$

    \If{$\|s_{\text{sit}} - s^{\text{eff}}_{\text{sit}}\|_2 > \theta_{\text{spike}}$}
        \State Emit situational \textsc{Spike} event $E_{\text{spike}}$
    \EndIf

    \If{$c_t \geq \gamma_{\text{short}}$ \textbf{and} $\bigl|\{e_i \in W \mid e_i^\top \tilde{e}_t \geq \tau\}\bigr| \geq 2$}
        \State Update the short-term state $s_{\text{short}}$; emit a \textsc{Shift} event (if $\ell_t \neq \ell_{t-1}$) or a \textsc{Reinforcement} event
    \EndIf

    \If{$c_t \geq \gamma_{\text{long}}$ \textbf{and} $\bigl|\{(e_i, c_i) \in W \mid e_i^\top \tilde{e}_t \geq \tau \;\wedge\; c_i \geq \gamma_{\text{long}}\}\bigr| \geq 3$}
        \State Update long-term state $s_{\text{long}}$; adapt baseline $b_{\text{long}} \gets b_{\text{long}} + \alpha\,(\tilde{e}_t - b_{\text{long}})$
        \State Emit a \textsc{Promotion} event
    \EndIf

    \State Update rolling window: $W \gets W \cup \{(\tilde{e}_t,\, c_t)\}$

    \For{each eligible event $E \in \{\textsc{Shift}, \textsc{Reinforcement}, \textsc{Promotion}\}$ with $\Delta_t^{\text{mag}} \geq 0.35$ and $c_t \geq 0.55$}
        \State Compute importance $w_t \gets \text{clip}_{[0,1]}(0.4 + 0.4\,\Delta_t^{\text{mag}} + 0.2\,c_t)$
        \State Write memory record $m$ into $\mathcal{M}$
    \EndFor

    \State \label{line:retrieval}Set the retrieval query $q \gets u_t$ (or use an explicit query, if provided)
    \State Query the memory store $\mathcal{M}$ and re-rank candidates using the hybrid score
    \Statex \qquad $S(m \mid q) = \dfrac{w_s\,\text{sim}(q, m) + w_r\,R(m) + w_i\,I(m)}{w_s + w_r + w_i}$
    \State Project the dominant label $\ell_t \gets \arg\max_l\, \hat{e}_t^\top \hat{q}_l$ and compute intensity $I_t \gets \text{clip}_{[0,1]}\!\left(\dfrac{\|e_t\|_2}{\sqrt{3}}\right)$
    \State Render the affective block $C_{\text{aff}}$ and assemble snapshot $\mathcal{S}$ within budget $B$
    \Statex \qquad $\mathcal{S} = \langle \text{Instructions},\, C_{\text{aff}},\, \mathcal{M}_{\text{top}},\, H_{\text{recent}} \rangle$
    \State \Return the context snapshot $\mathcal{S}$ and diagnostics $\mathcal{D}$
    \Statex \qquad $\mathcal{D} = \langle \text{token\_usage},\,
    |\mathcal{M}_{\text{top}}|,\, N_{\text{dropped}},\, \ell_t,\, I_t,\,
    \text{trend},\, c_t \rangle$
\end{algorithmic}
\end{breakablealgorithm}



The orchestration avoids performing affect processing and potentially slow classifier or memory operations while the main context-history lock is held. This design reduces the unnecessary serialization of unrelated context operations, while the emotion manager retains a separate lock for its internal state.



\subsection{Multi-label Emotion Classifier}

Let an input utterance be represented as

\begin{equation}
x = (w_1,w_2,\ldots,w_n)
\end{equation}

where $w_i$ denotes the $i$th token.

The objective is to predict the emotion vector

\begin{equation}
\mathbf{y}=[y_1,y_2,\ldots,y_{28}]
\end{equation}

where

\begin{equation}
y_i \in \{0,1\}
\end{equation}

indicates whether emotion $i$ is associated with the input utterance.

The model produces the logits

\begin{equation}
\mathbf{z}=f_{\theta}(x)
\end{equation}

The probability of each emotion is then calculated independently using the sigmoid function:

\begin{equation}
p_i = \frac{1}{1+e^{-z_i}}
\end{equation}

Emotion $i$ is selected when

\begin{equation}
p_i \geq \tau_i
\end{equation}

where $\tau_i$ represents the decision threshold for emotion $i$.

Unlike single-label classification, this formulation allows multiple emotion probabilities to exceed their corresponding thresholds simultaneously~\cite{demszky2020goemotions,huang2021seq2emo}.

The emotion-classification module uses DistilRoBERTa encoder~\cite{vaswani2017attention}, ~\cite{liu2019roberta}. Unlike a softmax classifier, which applies a sigmoid function independently to each logit and calculates a probability for every emotion:

\begin{equation}
P(E_i|x) = \sigma(z_i)
\end{equation}

where $\sigma$ denotes the sigmoid function.


Binary cross-entropy with logits is used for multi-label classification. The loss is defined as

\begin{equation}
L =
-\frac{1}{N}
\sum_{i=1}^{N}
\sum_{j=1}^{28}
[
y_{ij}\log(p_{ij})
+
(1-y_{ij})\log(1-p_{ij})
]
\end{equation}

where $N$ is the number of training examples, $y_{ij}$ is the ground-truth label for emotion $j$ in example $i$, and $p_{ij}$ is the corresponding predicted probability.

\subsection{Emotion Representation, Projection and Intensity}
\label{subsec:emotion_representation}


Classifier confidence is computed from the mean strength of the top-$n$
sigmoid probabilities and their margin over the next-highest probability.
Let the probabilities be sorted such that
$p_{t,(1)} \geq p_{t,(2)} \geq \cdots \geq p_{t,(L)}$, and define the mean
of the top-$n$ probabilities as
\begin{equation}
\bar{p}_{t,n}=\frac{1}{n}\sum_{k=1}^{n}p_{t,(k)}.
\label{eq:top_n_mean}
\end{equation}
The generalized confidence score is
\begin{equation}
c_t^{(n)} =
\operatorname{clip}_{[0,1]}
\left(
\frac{\bar{p}_{t,n}+
\left(\bar{p}_{t,n}-p_{t,(n+1)}\right)}{2}
\right).
\label{eq:confidence}
\end{equation}
Here, $1 \leq n < L$, where $L$ is the number of emotion labels. For
$n=1$, the expression reduces to the original top-1 confidence score based
on $p_{t,(1)}$ and its margin over $p_{t,(2)}$. This margin-aware measure
acts as a gate and modulates the magnitude of state updates.


Each classified conversational context is represented by an \texttt{EmotionSignal} that contains a VAD vector, classifier confidence, raw multi-label probabilities, and provenance information~\cite{mehrabian1996pad,russell1980circumplex}. The VAD vector is defined as
\begin{equation}
\mathbf{e}_t = [v_t,a_t,d_t]^\top,
\qquad
v_t,a_t,d_t\in[-1,1],
\end{equation}
where $v_t$ denotes valence, $a_t$ denotes arousal, and $d_t$ denotes dominance.

The raw probabilities are mapped into VAD space using an emotion-to-VAD prototype lexicon. Let $p_{t,l}$ denote the sigmoid probability of emotion label $l$ at turn $t$, and let $\mathbf{q}_l=[q_{l,v},q_{l,a},q_{l,d}]^\top$ denote its VAD prototype. The aggregated affect vector is
\begin{equation}
\tilde{\mathbf{e}}_t
=
\frac{\sum_{l}p_{t,l}\mathbf{q}_l}
{\sqrt{\max(1,\sum_l p_{t,l})}}.
\label{eq:probs_vad}
\end{equation}



For a non-neutral vector, dominant categorical emotion (without conflating affective direction with intensity),
\begin{equation}
\hat{\mathbf{e}}_t =
\frac{\mathbf{e}_t}{\|\mathbf{e}_t\|_2},
\qquad
\ell_t =
\arg\max_{l}
\left(
\hat{\mathbf{e}}_t^\top \hat{\mathbf{q}}_l
\right).
\label{eq:cosine_projection}
\end{equation}
The system returns a neutral label when the VAD magnitude is below the implementation's neutral floor of $0.12$. Intensity is calculated independently as
\begin{equation}
I_t =
\operatorname{clip}_{[0,1]}
\left(
\frac{\|\mathbf{e}_t\|_2}{\sqrt{3}}
\right),
\label{eq:intensity}
\end{equation}
where $\sqrt{3}$ is the maximum Euclidean magnitude of a VAD vector whose three components lie in $[-1,1]$.

\subsection{Tiered Temporal Affect Model}
\label{subsec:tiered_affect}

The affect tracker maintains three temporal tiers:
\begin{enumerate}
    \item \textbf{Situational affect}: captures immediate changes associated with the current interaction.
    \item \textbf{Short-term affect}: captures sustained recent emotional tendencies.
    \item \textbf{Long-term affect}: captures durable affective tendencies and is anchored to a slowly adapting baseline.
\end{enumerate}

Each tier stores a VAD state, a timestamp, a half-life, an anchor vector, and an independent multi-label probability-bar vector. In the supplied configuration, the default half-lives are $300$ seconds for the situational tier, $2{,}700$ seconds for the short-term tier, and $259{,}200$ seconds (three days) for the long-term tier. The corresponding learning rates are $0.70$, $0.35$, and $0.12$, respectively.



\subsubsection{Multi-Label Bar Vector} The multi-label \emph{bar vector} is a dictionary that maps each of the 28
GoEmotions-derived labels~\cite{demszky2020goemotions} to an independent intensity in $[0,100]$.
It is decayed and blended according to the same half-life and learning-rate schedule
as the VAD state of its tier, but it is never collapsed into a single value. The vector is
populated directly from the upstream classifier's raw multi-label
probabilities and is rendered in both the live prompt injection and persisted long-term memory text.


\subsubsection{Temporal Decay}
A stored VAD state decays exponentially toward its anchor. If $\mathbf{s}$ denotes the stored state, $\mathbf{b}$ its anchor, $h$ the tier-specific half-life, and $\Delta t$ the elapsed time, the effective state is
\begin{equation}
\mathbf{s}_{\mathrm{eff}}(t)
=
\mathbf{b}
+
\left(\mathbf{s}-\mathbf{b}\right)
2^{-\Delta t/h}.
\label{eq:vad_decay}
\end{equation}
The same half-life is applied to the multi-label bar vector, allowing the displayed probabilities to persist over time while remaining independent across labels.

\subsubsection{Confidence-Weighted State and Baseline Updates}
Both the tier state and the long-term baseline are updated through the same
exponential-blending rule. Let $\mathbf{x}_t$ denote the vector being updated
and let $\lambda_t$ denote its update coefficient. The unified update is
\begin{equation}
\mathbf{x}_{t}
=
\mathbf{x}_{t^-}
+
\lambda_t
\left(
\mathbf{e}_t-\mathbf{x}_{t^-}
\right),
\label{eq:state_update}
\end{equation}
where $\mathbf{x}_t=\mathbf{s}_t$ and
$\lambda_t=\eta_t$ for a tier-state update. The effective tier learning rate is
\begin{equation}
\eta_t = \eta_{\mathrm{tier}}c_t.
\label{eq:effective_lr}
\end{equation}
High-confidence observations therefore produce larger updates, whereas
uncertain observations have less influence. For long-term baseline
adaptation, $\mathbf{x}_t=\mathbf{b}_t$ and $\lambda_t=\alpha$. The same
tier-specific learning-rate and decay schedules are applied to the parallel
multi-label bar vector.

\subsubsection{Sustained-Support and Long-Term Promotion}
\label{subsec:support}

The architecture applies confidence gates rather than immediately promoting
every classifier prediction to a persistent affective state. The short-term
gate is the primary configurable value, with a default of
$\gamma_{\mathrm{short}}=0.45$. The situational and long-term gates are derived
from it as
\begin{align}
\gamma_{\mathrm{sit}}
&=\max\!\left(0.05,\gamma_{\mathrm{short}}-0.30\right)=0.15,\\
\gamma_{\mathrm{long}}
&=\min\!\left(0.95,\gamma_{\mathrm{short}}+0.15\right)=0.60.
\label{eq:confidence_gates}
\end{align}
Thus, situational, short-term, and long-term updates require confidence values
of at least $0.15$, $0.45$, and $0.60$, respectively.

For a short-term update, the current VAD signal must also be supported by two
prior observations. The implementation retains a rolling window of 12
VAD--confidence pairs. A prior observation $\mathbf{e}_i$ is considered
directionally consistent with the current signal $\mathbf{e}_t$ when
\begin{equation}
\mathbf{e}_i^\top\mathbf{e}_t \geq \tau,
\label{eq:agreement}
\end{equation}


where the default agreement threshold is $\tau=0.25$. Therefore, the default
short-term support requirement is $k_{\mathrm{stm}}=2$. The current observation
is appended to the rolling window only after the support check, preventing it
from satisfying its own corroboration requirement.

Long-term promotion requires $k_{\mathrm{ltm}}=3$ prior high-confidence observations that
satisfy the same directional-agreement condition. Following promotion, the
state is updated with the smaller long-term learning rate. The persistent
baseline is adapted using the unified rule in Eq.~\ref{eq:state_update}, with
$\mathbf{x}_t=\mathbf{b}_t$ and $\lambda_t=\alpha$. The default value is
$\alpha=0.03$. The long-term tier subsequently decays toward this evolving
baseline rather than toward the neutral origin. 



\subsection{Event Detection and Emotion-Aware Memory Writing}
\label{sec:method-events}

A situational \emph{Spike} is emitted when the state-change magnitude exceeds 0.25; short-term events are classified as \emph{Shift} or \emph{Reinforcement} depending on whether the projected label changes; a successful long-term update emits a \emph{Promotion}. Eligible events (short-term reinforcement/shift, long-term promotion -- one-off situational spikes are excluded) are written to long-term memory only when the change magnitude is $\geq 0.35$ and confidence $\geq 0.55$, subject to a five-minute per-tier/label cooldown. Memory importance is

\begin{equation}
w_t = \mathrm{clip}_{[0,1]}\!\left(0.4 + 0.4\,\Delta_t + 0.2\,c_t\right)\;.
\label{eq:importance}
\end{equation}









% =========================================================
% IV. EXPERIMENTAL RESULTS AND EVALUATION
% =========================================================
\section{Experimental Results and Evaluation}

The reference implementation is written in Python with no required
third-party dependencies for its core logic (a hashing-based embedder and
heuristic token counter). The
classifier itself is injected via a narrow protocol
(\texttt{classify(turns) -> EmotionSignal}) carrying both a VAD vector
and the raw multi-label probability dictionary, so that a
fine-tuned DistilRoBERTa model on
GoEmotions~\cite{sanh2019distilbert,liu2019roberta,demszky2020goemotions}
can be attached without modifying the memory architecture. 

\subsection{Multi-label Emotion Clssifier model}

\subsubsection{Dataset}

The classifier model uses GoEmotions as the primary training dataset rather than single-label annotation; binary-sentiment corpora and contextual/multimodal corpora (EmotionLines, MELD, MEISD \cite{chen2018emotionlines,poria2019meld,zhang2020meisd}) do not provide the required broad multi-label taxonomy. GoEmotions contains approximately 58,000 English Reddit comments annotated using 27 emotion categories and a neutral category \~cite{demszky2020goemotions}.

The complete label space contains 28 labels:

\begin{center}
\begin{tabular}{@{}p{0.44\columnwidth}p{0.44\columnwidth}@{}}
\textbullet\ admiration     & \textbullet\ amusement \\
\textbullet\ anger          & \textbullet\ annoyance \\
\textbullet\ approval       & \textbullet\ caring \\
\textbullet\ confusion      & \textbullet\ curiosity \\
\textbullet\ desire         & \textbullet\ disappointment \\
\textbullet\ disapproval    & \textbullet\ disgust \\
\textbullet\ embarrassment & \textbullet\ excitement \\
\textbullet\ fear           & \textbullet\ gratitude \\
\textbullet\ grief          & \textbullet\ joy \\
\textbullet\ love           & \textbullet\ nervousness \\
\textbullet\ optimism       & \textbullet\ pride \\
\textbullet\ realization    & \textbullet\ relief \\
\textbullet\ remorse        & \textbullet\ sadness \\
\textbullet\ surprise       & \textbullet\ neutral
\end{tabular}
\end{center}



\subsubsection{LoRA Fine-Tuning}

Low-Rank Adaptation is used to adapt the pretrained emotion model without performing full-parameter fine-tuning \cite{hu2021lora}.

LoRA represents the weight update as:

\begin{equation}
W' = W + BA
\end{equation}


Here, $B$ and $A$ are low-rank trainable matrices.

\begin{table}[h]
\centering
\caption{LoRA configuration and training hyperparameters.}
\label{tab:lora-config}
\footnotesize
\begin{tabular}{@{}>{\raggedright\arraybackslash}p{0.32\columnwidth}
                    >{\raggedright\arraybackslash}p{0.62\columnwidth}@{}}
\toprule
\textbf{Hyperparameter} & \textbf{Value} \\
\midrule
Base checkpoint & \url{j-hartmann/emotion-english-distilroberta-base} \\
LoRA rank ($r$) & 16 \\
LoRA alpha & 32 \\
LoRA dropout & 0.1 \\
Target modules & query, key, value \\
Fully-trained modules & classifier head (\texttt{modules\_to\_save}) \\
Trainable parameters & 1,054,492 / 83,194,424 (1.27\%) \\
Epochs & 4 \\
Per-device batch size & 32 (T4 $\times$ 2) \\
Learning rate & $2 \times 10^{-4}$ \\
Weight decay & 0.01 \\
Warmup ratio & 0.06 \\
Precision & fp16 \\
Loss function & Binary cross-entropy with logits (per label) \\
\bottomrule
\end{tabular}
\end{table}

The original pretrained parameters remain frozen while the LoRA parameters are optimized during training.

The classification head is also trained because the original classification head does not correspond to the proposed 28-label output space.



\subsubsection{Fine-tuned Model Evaluation Results}

The model was evaluated on the official GoEmotions test split containing 5,427 examples~\cite{demszky2020goemotions}.

The results are presented in Table~\ref{tab:results}.

\begin{table}[h]
\centering
\caption{Test-set performance at a 0.5 sigmoid threshold, $n = 5{,}427$.}
\label{tab:results}
\begin{tabular}{@{}lll@{}}
\toprule
\textbf{Metric} & \textbf{Value} & \textbf{Interpretation} \\
\midrule
F1 (micro) & 0.5584 & Aggregate correctness across all label decisions \\
F1 (macro) & 0.3859 & \begin{tabular}[t]{@{}l@{}}Unweighted average across the 28 labels;\\ sensitive to rare-class performance\end{tabular} \\
Precision (micro) & 0.7151 & Of predicted labels, 71.5\% were correct \\
Recall (micro) & 0.4581 & Of true labels, 45.8\% were recovered at threshold 0.5 \\
Hamming loss & 0.0302 & \begin{tabular}[t]{@{}l@{}}3.0\% of individual label decisions (28 per\\ example) were wrong\end{tabular} \\
Exact match & 0.4293 & \begin{tabular}[t]{@{}l@{}}42.9\% of test examples had their entire label\\ set predicted exactly right\end{tabular} \\
\bottomrule
\end{tabular}
\end{table}


The micro-precision of 0.7151 indicates that the model is relatively conservative when predicting emotions. However, the micro-recall of 0.4581 indicates that a considerable number of relevant emotions are not predicted at the selected threshold.

The difference between micro-F1 and macro-F1 also indicates the influence of class imbalance. Frequent emotions contribute more strongly to micro-level performance, whereas rare emotions are more difficult to classify.


\subsection{Short-Term Confidence-Gate Admission}
The confidence-gate sensitivity analysis confirms that the production default
$\gamma_{\mathrm{short}}=0.45$ yields a mean short-term admission rate of
$0.1525\pm0.0773$, with a 95\% confidence interval of
$[0.1186,\,0.1864]$. This value quantifies the proportion of candidate signals
admitted to short-term affect updating under the default gate; it is distinct
from the two-observation sustained-support requirement, which is applied as an
additional corroboration condition.


\subsection{Decay Correctness}
The tiered design depends on decay being composable: decaying by elapsed
time $t_1+t_2$ in one update must equal decaying by $t_1$ then $t_2$ in
two successive updates, since tiers are updated at irregular intervals
determined by when messages arrive, not on a fixed clock. We verified
this semigroup property across 24 (half-life, elapsed-time)
configurations spanning all three tiers' default half-lives. The maximum
absolute error observed was $2.5\times10^{-16}$ (floating-point noise
floor), confirming the decay implementation is a correct, order-
independent exponential process.

\subsection{Label-Projection Accuracy: Defect and Correction}
An initial implementation matched a VAD point to its nearest emotion
prototype by raw Euclidean distance. Because prototype vectors have
different magnitudes, this systematically biases classification toward
small-magnitude prototypes regardless of direction. We quantified this
with a controlled test: for each of 8 non-neutral prototype emotions, we
generated points along that prototype's direction (Gaussian angular
jitter, $\sigma=0.05$) at 20 magnitudes from 0.05 to 1.0, 30 samples per
cell ($n=4{,}800$), and measured how often each projector recovered the
generating prototype's label.

\begin{table}[h]
\caption{Label projection accuracy: original vs. corrected}
\centering
\begin{tabular}{lcc}
\toprule
Regime & Original (Euclidean NN) & Corrected (cosine) \\
\midrule
Overall ($n=4{,}800$) & 44.8\% & \textbf{90.0\%} \\
Weak signal ($\|v\|\leq0.3$) & 0.0\% & \textbf{66.7\%} \\
Strong signal ($\|v\|>0.7$) & 90.9\% & \textbf{100.0\%} \\
\bottomrule
\end{tabular}
\end{table}

The fix separates direction from magnitude: the corrected projector
scores label direction by cosine similarity against unit-normalized
prototypes, and reports the point's own magnitude as intensity, rather
than conflating the two. Figure~1 shows accuracy as a function of signal
magnitude; the original projector is at floor accuracy below
$\|v\|\approx0.3$ --- precisely the regime a single, moderately-confident
conversational turn is likely to occupy --- and only approaches ceiling
accuracy for unusually intense signals.

\begin{figure}[h]
\centering
\includegraphics[width=0.95\linewidth]{fig1_label_projection.png}
\caption{Label projection accuracy vs.\ signal magnitude, original
(Euclidean nearest-neighbor) vs.\ corrected (cosine similarity)
projector. $n=4{,}800$ synthetic points, ground truth by construction.}
\label{fig:labelproj}
\end{figure}

\subsection{Tiered Settling Dynamics}
We confirmed the intended separation of timescales by injecting a
constant, high-confidence signal every 60 simulated seconds (via a
controllable clock) and recording the turn and elapsed time at which
each tier's magnitude first reached 90\% of the signal's own magnitude.

\begin{table}[h]
\caption{Turns and simulated time to reach 90\% of asymptotic signal}
\centering
\begin{tabular}{lcc}
\toprule
Tier & Turns to 90\% & Simulated time \\
\midrule
Situational & 2 & 120 s \\
Short-term & 8 & 480 s (8 min) \\
Long-term & 21 & 1{,}260 s (21 min) \\
\bottomrule
\end{tabular}
\end{table}

The ordering situational $<$ short-term $<$ long-term held at every
configuration tested. Figure~2 additionally shows an emergent property
not directly specified by the design: situational and short-term settle
to a dynamic equilibrium slightly below the asymptote ($\approx$93--97\%),
because inter-turn decay is non-negligible relative to the per-turn
blend step at a 60-second turn interval, whereas long-term eventually
exceeds them, because its decay anchor (the trait baseline) itself
migrates toward the sustained signal over time.

\begin{figure}[h]
\centering
\includegraphics[width=0.95\linewidth]{fig2_tier_dynamics.png}
\caption{Tier settling dynamics under a sustained signal, injected every
60 simulated seconds via a controllable clock.}
\label{fig:tierdynamics}
\end{figure}



% =========================================================
% V. DISCUSSION
% =========================================================
\section{Discussion}

\subsection{Failure Analysis of the Fine-Tuned Emotion Classifier}

Despite reasonable aggregate performance, the fine-tuned classifier exhibits three interrelated failure modes that limit its practical utility. The large macro-F1 (0.386) versus micro-F1 (0.558) gap reflects systematic under-detection of rare labels (e.g., grief, remorse), which the unweighted cross-entropy objective sacrifices in favor of frequent classes like neutral and annoyance, causing rare emotions to be misattributed or missed entirely. This is compounded by the default 0.5 sigmoid threshold, where micro-precision (0.715) far exceeds micro-recall (0.458), silently dropping legitimate co-occurring emotions whose probabilities fall below cutoff—directly undermining the project's goal of surfacing mixed emotional signals. Finally, a nontrivial share of apparent errors likely stems from genuine taxonomic ambiguity rather than misclassification, as GoEmotions' own low inter-annotator agreement ($\kappa \approx 0.29$--$0.34$) and semantically overlapping label pairs (e.g., annoyance--anger--disapproval) introduce boundary noise that exact-match metrics cannot distinguish from true failures. 


\subsection{Failure Modes of Embedding-Based Long-Term Affective Memory}
Once an emotional episode is written to the long-term store, it is
represented, embedded, and retrieved identically to an ordinary fact,
which introduces four specific problems. First, retrieval is driven by
text/topic similarity, not by affective similarity: the stored VAD point
is not retained in memory metadata by default, so a query intended to
find ``times the user felt anxious'' competes only on weak lexical
grounds against topically similar but affectively different content.
This is straightforward to address by storing the raw VAD alongside the
text and adding an explicit affect-similarity term to the retrieval
score. Second, the tiered half-life structure that gives situational,
short-term, and long-term affect their distinct temporal meaning is
discarded on write; the generic recency term used by the retrieval score
treats an emotional episode exactly like a static fact, decaying it at
one global rate rather than a rate appropriate to how quickly affective
relevance actually fades. Third, only discrete episodes are ever written
--- the aggregate trait baseline, which is the closest representation to
``who this person generally is, emotionally,'' is never itself expressed
as retrievable text, so semantic search can recover individual incidents
but not trends. Fourth, because generic text similarity dominates the
hybrid retrieval score, an emotionally significant memory can be
outranked by an unrelated but more lexically similar fact unless a
dedicated retrieval budget is reserved for affective content, mirroring
the fixed budget already reserved for the live affect block
(Section~III-C).

\subsection{Context Management Efficiency}
Maintaining emotional state alongside conversational context does not
require growing the token budget in proportion to conversation length.
In this architecture, only two things ever enter the model's prompt on
account of affect: a small, fixed-size live summary of the current
situational and short-term state, capped at a fixed fraction of the
total budget regardless of session length, and, when relevant, a handful
of retrieved long-term emotional episodes competing for space in the
same token-bounded memory block as any other fact. This keeps the
\emph{marginal} cost of affect-awareness constant per turn rather than
cumulative: Experiment~5.4 measured context-assembly latency scaling
linearly with the number of stored long-term memories, with 100\%
token-budget compliance across all tested scales, and this cost is
identical whether or not any given memory happens to be emotional in
origin, since all memories share one retrieval pipeline. The efficiency
tradeoff is therefore paid not in context-window size but in write-path
selectivity: the confidence, magnitude, and cooldown thresholds governing
the memory writer determine how much of the emotional signal computed
every turn ever becomes a persisted, retrievable memory, versus remaining
a transient contribution to the decaying in-memory tier state. This
separation of cheap, bounded, always-present live affect from selective,
embedding-indexed episodic affect is what allows the architecture to
sustain emotional continuity across arbitrarily long sessions without
context-management cost growing with session length.



\section{Limitations}

Several limitations bound the claims made in this paper. First, all
reported experiments use synthetic inputs and a rule-based classifier
stand-in rather than a trained model; they validate the memory and decay
architecture, not the classification accuracy of any specific emotion
model. Standard classifier evaluation --- micro/macro F1, per-label
precision and recall, and calibration analysis against a held-out
GoEmotions~\cite{demszky2020goemotions} split --- is necessary before any claim about real-world emotion-recognition accuracy can be made, and is left as future work. 

Second, no experiment in this paper involves real
users or real conversational data; any claim about ecological validity,
response quality, or user experience requires a human evaluation or
annotated real-conversation corpus, which we do not report. Third,
results are reported from single-seed runs; we do not report confidence
intervals or significance tests, and recommend multi-seed replication
with bootstrap intervals before treating reported deltas (e.g.,
$44.8\%\rightarrow90.0\%$) as more than descriptive. 

Fourth, the corrected cosine-based label projector still exhibits a residual failure
mode near tier decision boundaries: during a rapid mood transition, a
blended vector average of two emotions can transiently project onto a
third, semantically unrelated label (Appendix~A, Scenario F); the
independent multi-label bar representation is unaffected by this and
should be preferred over the single projected label wherever the two
disagree. 

Fifth, inferred emotional state persisted to long-term memory
is sensitive data; no redaction, consent, or differential retention
policy is implemented for affective memories specifically, and the
long-term store has no forgetting or pruning mechanism, so it grows
without bound over the lifetime of a deployment.


Sixth, A single global threshold may not be optimal for all emotion categories, particularly when label frequencies and decision boundaries differ~\cite{huang2021limits,huang2021seq2emo}.
    

% =========================================================
% VI. CONCLUSION
% =========================================================


\section{Conclusion}
\label{sec:conclusion}

This work presented an end-to-end emotion-aware conversational system integrating a LoRA fine-tuned DistilRoBERTa multi-label classifier with a tiered affective memory architecture combining a continuous VAD reasoning core with independently decaying multi-label representations across three timescales. Controlled experiments validated the memory module's decay correctness, timescale separation, and linear scalability while identifying and correcting a magnitude-biased label-projection defect; the classifier's precision--recall asymmetry and macro/micro-F1 gap point to threshold calibration and class reweighting as immediate priorities. Future work should address these classifier-level limitations alongside conversational emotion trajectories, stronger safety detection, and systematic human evaluation of end-to-end response quality.


% =========================================================
% REFERENCES
% =========================================================

\begin{thebibliography}{99}

\bibitem{lewis2020rag}
P. Lewis \textit{et al.}, ``Retrieval-Augmented Generation for
Knowledge-Intensive NLP Tasks,'' in \textit{Advances in Neural Information
Processing Systems}, vol. 33, 2020.

\bibitem{packer2023memgpt}
C. Packer, S. Wooders, K. Lin, V. Fang, S. G. Patil, I. Stoica, and
J. E. Gonzalez, ``MemGPT: Towards LLMs as Operating Systems,''
\textit{arXiv preprint arXiv:2310.08560}, 2023.

\bibitem{park2023generative}
J. S. Park, J. C. O'Brien, C. J. Cai, M. R. Morris, P. Liang, and
M. S. Bernstein, ``Generative Agents: Interactive Simulacra of Human
Behavior,'' in \textit{Proc. 36th Annual ACM Symposium on User Interface
Software and Technology (UIST)}, 2023, pp. 1--22.

\bibitem{rashkin2019empathetic}
H. Rashkin, E. M. Smith, M. Li, and Y.-L. Boureau, ``Towards Empathetic
Open-Domain Conversation Models: A New Benchmark and Dataset,'' in
\textit{Proc. 57th Annual Meeting of the Association for Computational
Linguistics (ACL)}, 2019, pp. 5370--5381,
doi: 10.18653/v1/P19-1534.

\bibitem{liu2021esconv}
S. Liu, C. Zheng, O. Demasi, S. Sabour, Y. Li, Z. Yu, Y. Jiang, and
M. Huang, ``Towards Emotional Support Dialog Systems,'' in
\textit{Proc. 59th Annual Meeting of the Association for Computational
Linguistics and the 11th International Joint Conference on Natural
Language Processing}, 2021, pp. 3469--3483,
doi: 10.18653/v1/2021.acl-long.269.

\bibitem{zhang2020meisd}
S. Zhang \textit{et al.}, ``MEISD: A Multimodal Multi-Label Emotion,
Intensity and Sentiment Dialogue Dataset for Emotion Recognition and
Sentiment Analysis in Conversations,'' in \textit{Proc. 28th International
Conference on Computational Linguistics (COLING)}, 2020.

\bibitem{huang2021seq2emo}
C. Huang, A. Trabelsi, X. Qin, N. Farruque, L. Mou, and O. Zaiane,
``Seq2Emo: A Sequence to Multi-Label Emotion Classification Model,'' in
\textit{Proc. 2021 Conference of the North American Chapter of the
Association for Computational Linguistics: Human Language Technologies},
2021.

\bibitem{demszky2020goemotions}
D. Demszky, D. Movshovitz-Attias, J. Ko, A. Cowen, G. Nemade, and
S. Ravi, ``GoEmotions: A Dataset of Fine-Grained Emotions,'' in
\textit{Proc. 58th Annual Meeting of the Association for Computational
Linguistics (ACL)}, 2020, pp. 4040--4054.

\bibitem{wu2025longmemeval}
D. Wu, H. Wang, W. Yu, Y. Zhang, K.-W. Chang, and D. Yu,
``LongMemEval: Benchmarking Chat Assistants on Long-Term Interactive
Memory,'' in \textit{Proc. International Conference on Learning
Representations (ICLR)}, 2025.

\bibitem{edge2024graphrag}
D. Edge, H. Trinh, N. Cheng, J. Bradley, A. Chao, A. Mody, S. Truitt,
and J. Larson, ``From Local to Global: A Graph RAG Approach to
Query-Focused Summarization,'' \textit{arXiv preprint arXiv:2404.16130},
2024.

\bibitem{rasmussen2025zep}
P. Rasmussen, P. Paliychuk, T. Beauvais, J. Ryan, and D. Chalef,
``Zep: A Temporal Knowledge Graph Architecture for Agent Memory,''
\textit{arXiv preprint arXiv:2501.13956}, 2025.

\bibitem{majumder2019dialoguernn}
N. Majumder, S. Poria, D. Hazarika, R. Mihalcea, A. Gelbukh, and
E. Cambria, ``DialogueRNN: An Attentive RNN for Emotion Detection in
Conversations,'' \textit{Proc. AAAI Conference on Artificial Intelligence},
vol. 33, no. 1, pp. 6818--6825, 2019,
doi: 10.1609/aaai.v33i01.33016818.

\bibitem{ghosal2020cosmic}
D. Ghosal, N. Majumder, A. Gelbukh, R. Mihalcea, and S. Poria,
``COSMIC: COmmonSense Knowledge for eMotion Identification in
Conversations,'' in \textit{Findings of the Association for Computational
Linguistics: EMNLP 2020}, 2020, pp. 2470--2481,
doi: 10.18653/v1/2020.findings-emnlp.224.

\bibitem{liu2019roberta}
Y. Liu \textit{et al.}, ``RoBERTa: A Robustly Optimized BERT Pretraining
Approach,'' \textit{arXiv preprint arXiv:1907.11692}, 2019.

\bibitem{hu2021lora}
E. J. Hu, Y. Shen, P. Wallis, Z. Allen-Zhu, Y. Li, S. Wang, L. Wang,
and W. Chen, ``LoRA: Low-Rank Adaptation of Large Language Models,''
\textit{arXiv preprint arXiv:2106.09685}, 2021.

\bibitem{russell1980circumplex}
J. A. Russell, ``A Circumplex Model of Affect,'' \textit{Journal of
Personality and Social Psychology}, vol. 39, no. 6, pp. 1161--1178, 1980.

\bibitem{mehrabian1996pad}
A. Mehrabian, ``Pleasure-Arousal-Dominance: A General Framework for
Describing and Measuring Individual Differences in Temperament,''
\textit{Current Psychology}, vol. 14, no. 4, pp. 261--292, 1996.

\bibitem{zhong2024memorybank}
W. Zhong, L. Guo, Q. Gao, H. Ye, and Y. Wang, ``MemoryBank: Enhancing
Large Language Models with Long-Term Memory,'' \textit{Proc. AAAI
Conference on Artificial Intelligence}, vol. 38, no. 17,
pp. 19724--19731, 2024, doi: 10.1609/aaai.v38i17.29946.

\bibitem{Fathi2026}
A. Fathi, S. Nouri, and S. Danesh, ``Adaptive Memory Systems for
Personalized Chatbots,'' \textit{International Journal of Computational
Health and Machine Learning}, vol. 4, no. 1, 2026.

\bibitem{Oni2025}
O. Oni, ``Memory-Enhanced Conversational AI: A Generative Approach for
Context-Aware and Personalized Chatbots,'' \textit{Communication in
Physical Sciences}, vol. 12, no. 2, pp. 649--657, 2025.

\bibitem{Jover2026}
S. M. Jover, ``Empathic Conversational AI: A Voice Assistant with Emotion
Awareness and Persistent Memory for Elderly People with Cognitive
Decline,'' Bachelor's thesis, LAB University of Applied Sciences, 2026.

\bibitem{Jiang2026}
Z. Jiang, C. Zheng, H. Chen, and L. Chen, ``RECALLbot: Designing Agentic
Memory and Reciprocal Disclosure for Human-Chatbot Relationships,'' in
\textit{Proc. 2026 CHI Conference on Human Factors in Computing Systems},
2026.

\bibitem{Cox2026}
S. R. Cox, M. Yin, and N. van Berkel, ``How Does Chatbot Memory and
Framing Affect People's Self-Disclosure Decisions?'' in \textit{Proc.
Bias4Trust Workshop at CHI 2026}, 2026.

\bibitem{gutierrez2024hipporag}
B. J. Guti\'errez, Y. Shu, Y. Gu, M. Yasunaga, and Y. Su, ``HippoRAG:
Neurobiologically Inspired Long-Term Memory for Large Language Models,''
\textit{arXiv preprint arXiv:2405.14831}, 2024.

\bibitem{xu2025amem}
W. Xu, Z. Liang, K. Mei, H. Gao, J. Tan, and Y. Zhang, ``A-MEM: Agentic
Memory for LLM Agents,'' \textit{arXiv preprint arXiv:2502.12110}, 2025.

\bibitem{chhikara2025mem0}
P. Chhikara, D. Khant, S. Aryan, T. Singh, and D. Yadav, ``Mem0:
Building Production-Ready AI Agents with Scalable Long-Term Memory,''
\textit{arXiv preprint arXiv:2504.19413}, 2025.

\bibitem{maharana2024longterm}
A. Maharana, D.-H. Lee, S. Tulyakov, M. Bansal, F. Barbieri, and
Y. Fang, ``Evaluating Very Long-Term Conversational Memory of LLM
Agents,'' in \textit{Proc. 62nd Annual Meeting of the Association for
Computational Linguistics (ACL)}, vol. 1, 2024, pp. 13851--13870.

\bibitem{Akram2026}
R. Akram, T. Ali, N. A. Khan, H. A. Khan, A. Hasnain, and K. Zahra,
``Enhancing Human--Computer Interaction Through Emotion Detection in
Chatbots,'' \textit{Saudi Journal of Engineering and Technology}, vol. 11,
no. 4, pp. 312--329, 2026, doi: 10.36348/sjet.2026.v11i04.016.

\bibitem{Abinaya2025}
S. Abinaya, K. S. Ashwin, and A. Sherly Alphonse, ``Enhanced
Emotion-Aware Conversational Agent: Analyzing User Behavioral Status for
Tailored Responses in Chatbot Interactions,'' \textit{IEEE Access},
vol. 13, pp. 19770--19787, 2025,
doi: 10.1109/ACCESS.2025.3534197.

\bibitem{Devi2026}
S. S. Devi, R. Lavanya, G. Santhiya, M. Malleswari, N. Soundariya, and
S. K. Manikandan, ``A Multimodal Context-Aware AI Conversational Agent
for Emotion Recognition and Adaptive Stress Management,'' in \textit{Proc.
2026 International Conference on Emerging Trends in Mobile Computing and
Sustainable Informatics (ICEMCSI)}, 2026,
doi: 10.1109/ICEMCSI67638.2026.11603222.

\bibitem{chen2018emotionlines}
S.-Y. Chen, C.-C. Hsu, C.-C. Kuo, T.-H. Huang, and L.-W. Ku,
``EmotionLines: An Emotion Corpus of Multi-Party Conversations,''
\textit{arXiv preprint arXiv:1802.08379}, 2018.

\bibitem{poria2019meld}
S. Poria, N. Majumder, R. Mihalcea, and E. Hovy, ``MELD: A Multimodal
Multi-Party Dataset for Emotion Recognition in Conversations,'' in
\textit{Proc. 57th Annual Meeting of the Association for Computational
Linguistics (ACL)}, 2019.

\bibitem{Raghav2026}
J. Raghav, N. G., and L. S. Raveendran, ``Design and Evaluation of an
Emotionally Intelligent Conversational Agent With Long-Term Context
Awareness,'' in \textit{Proc. 1st International Conference on Intelligent
Computing, Networks and Security (IC-ICNS)}, 2026,
doi: 10.1109/IC-ICNS68863.2026.11537630.

\bibitem{ChangHerath2025}
F. Chang and D. Herath, ``From Interaction to Relationship: The Role of
Memory, Learning, and Emotional Intelligence in AI-Embodied Human
Engagement,'' in \textit{Proc. 20th ACM/IEEE International Conference on
Human-Robot Interaction (HRI)}, Melbourne, Australia, 2025,
pp. 1269--1270, doi: 10.1109/HRI61500.2025.10973813.

\bibitem{vaswani2017attention}
A. Vaswani \textit{et al.}, ``Attention Is All You Need,'' in
\textit{Advances in Neural Information Processing Systems}, 2017.

\bibitem{sanh2019distilbert}
V. Sanh, L. Debut, J. Chaumond, and T. Wolf, ``DistilBERT, a Distilled
Version of BERT: Smaller, Faster, Cheaper and Lighter,'' \textit{arXiv
preprint arXiv:1910.01108}, 2019.

\bibitem{huang2021limits}
C. Huang, A. Trabelsi, X. Qin, N. Farruque, L. Mou, and O. Zaiane,
``Uncovering the Limits of Text-Based Emotion Detection,'' in
\textit{Findings of the Association for Computational Linguistics: EMNLP},
2021.

\end{thebibliography}

\end{document}
